\documentclass[oneside,a4paper,14pt]{extarticle}

\usepackage[T2A]{fontenc}
\usepackage[top=20mm,bottom=20mm,left=30mm,right=10mm]{geometry}
\usepackage[utf8]{inputenc}
\usepackage[russian]{babel}
\usepackage{textcomp}
\usepackage{indentfirst}
\usepackage{graphicx}
\usepackage{mwe}
\usepackage{wrapfig}
\usepackage{caption}
\usepackage{amsmath}
\usepackage{amsfonts}
\usepackage{amsthm}
\usepackage[all]{xy}
\usepackage[breaklinks,hidelinks]{hyperref}
\usepackage{titlesec}
\usepackage{listings}
\usepackage{minted,fancyvrb}
\usepackage{float}
\usepackage{array}
\usepackage{url}
\captionsetup{labelsep=endash,justification=centering}

\setminted{style=rainbow_dash,fontsize=\small}

\titleformat{\section}
{\normalsize\bfseries}
{\thesection}{1em}{}
\titleformat{\subsection}
{\normalsize\bfseries}
{\thesubsection}{1em}{}
\titleformat{\subsubsection}
{\normalsize\bfseries}
{\thesubsubsection}{1em}{}

\renewcommand\baselinestretch{1.33}\normalsize
\setlength{\parindent}{1.25cm}
\setlength{\emergencystretch}{3em}

\begin{document}
\newpage\thispagestyle{empty}

\begin{center}
МИНИСТЕРСТВО НАУКИ И ВЫСШЕГО ОБРАЗОВАНИЯ\\
РОССИЙСКОЙ ФЕДЕРАЦИИ\\
ФЕДЕРАЛЬНОЕ ГОСУДАРСТВЕННОЕ БЮДЖЕТНОЕ\\
ОБРАЗОВАТЕЛЬНОЕ УЧРЕЖДЕНИЕ ВЫСШЕГО ОБРАЗОВАНИЯ\\
«ВЯТСКИЙ ГОСУДАРСТВЕННЫЙ УНИВЕРСИТЕТ»\\
Институт математики и информационных систем\\
Факультет автоматики и вычислительной техники\\
Кафедра электронных вычислительных машин
\end{center}

\vspace{30mm}

\begin{center}
Отчёт по лабораторной работе № 2.2\\
по дисциплине\\
«Управление данными»\\
«Основы DML-запросов в PostgreSQL»
\end{center}

\vspace{40mm}

\noindent
\begin{tabular}{l r}
Выполнил студент гр. ИВТб-2303-06-00 \hspace{3mm}\rule[-1mm]{20mm}{0.15mm} /Перевозчиков Д.А./ \\

Проверил \hspace{57mm}\rule[-1mm]{20mm}{0.15mm} /Клюкин В. Л./
\end{tabular}

\vfill
\begin{center}
Киров\\
2026
\end{center}


\newpage\thispagestyle{empty}

\section{Цели лабораторной работы}

\begin{itemize}
    \item научиться использовать агрегатные функции;
    \item научиться использовать операторы и встроенные функции, работать с датами;
    \item научиться использовать подзапросы и предложения \texttt{WITH}, \texttt{RETURNING};
    \item рассмотреть и выполнить различные планы выполнения запросов;
    \item выполнить \texttt{FORMAT JSON} для выборки из нескольких связанных таблиц;
    \item выполнить три способа соединения для одной и той же выборки и определить наиболее оптимальный.
\end{itemize}

\section{Задание}

Выполнить все примеры из методических указаний для своей предметной области, используя данные и знания, полученные в предыдущих лабораторных работах. В качестве предметной области используется база данных гоночного автоклуба, созданная в лабораторной работе № 1 и заполненная данными в лабораторной работе № 2.1. База данных содержит таблицы \texttt{car}, \texttt{trainer}, \texttt{trainer\_car}, \texttt{competition} и \texttt{achievement}.

\newpage
\section{Решение}

\subsection{Агрегатные функции и предложение HAVING}

Агрегатные функции вычисляют одно значение по множеству входных строк. Для получения результата отдельно по каждой группе строк используется предложение \texttt{GROUP BY}. Условия, в которых участвует результат агрегатной функции, записываются в предложении \texttt{HAVING}, а не в \texttt{WHERE}: \texttt{WHERE} фильтрует отдельные строки до группировки, а \texttt{HAVING} --- уже сформированные группы.

\textbf{Запрос 1.} Подсчёт количества автомобилей в автопарке с помощью \texttt{COUNT(*)}.

\begin{minted}[breaklines]{sql}
SELECT COUNT(*) AS total_cars
FROM car;
\end{minted}

\begin{table}[H]
\centering
\footnotesize

\caption{Результат запроса 1}\label{tab:agg1}
\begin{tabular}{|l|}
\hline
\textbf{total\_\allowbreak cars} \\ \hline
12 \\ \hline
\end{tabular}
\end{table}

В автопарке числится 12 автомобилей, включая автомобили, отправленные в утилизацию.


\textbf{Запрос 2.} Вычисление среднего, минимального, максимального и суммарного значений мощности для автомобилей, не находящихся в утилизации.

\begin{minted}[breaklines]{sql}
SELECT
    AVG(horsepower) AS average_horsepower,
    MIN(horsepower) AS min_horsepower,
    MAX(horsepower) AS max_horsepower,
    SUM(horsepower) AS total_horsepower
FROM car
WHERE is_utilized = FALSE;
\end{minted}

\begin{table}[H]
\centering
\footnotesize

\caption{Результат запроса 2}\label{tab:agg2}
\begin{tabular}{|l|l|l|l|}
\hline
\textbf{average\_\allowbreak horsepower} & \textbf{min\_\allowbreak horsepower} & \textbf{max\_\allowbreak horsepower} & \textbf{total\_\allowbreak horsepower} \\ \hline
258.8181818181818182 & 76 & 510 & 2847 \\ \hline
\end{tabular}
\end{table}

Средняя мощность рабочих автомобилей составляет около 258,8 л.с., минимальная --- 76 л.с., максимальная --- 510 л.с.


\textbf{Запрос 3.} Группировка автомобилей по цвету. Для каждого цвета определяется количество автомобилей и их средняя мощность, результат сортируется по убыванию количества.

\begin{minted}[breaklines]{sql}
SELECT
    color,
    COUNT(*) AS car_count,
    AVG(horsepower) AS average_horsepower
FROM car
WHERE is_utilized = FALSE
GROUP BY color
ORDER BY car_count DESC;
\end{minted}

\begin{table}[H]
\centering
\footnotesize

\caption{Результат запроса 3}\label{tab:agg3}
\begin{tabular}{|l|l|l|}
\hline
\textbf{color} & \textbf{car\_\allowbreak count} & \textbf{average\_\allowbreak horsepower} \\ \hline
Красный & 3 & 164 \\ \hline
Оранжевый & 2 & 265 \\ \hline
Белый & 2 & 135 \\ \hline
Синий & 2 & 297.5 \\ \hline
Серый & 1 & 510 \\ \hline
Чёрный & 1 & 450 \\ \hline
\end{tabular}
\end{table}


\textbf{Запрос 4.} Группировка автомобилей по тренерам. Функция \texttt{COUNT} считает количество автомобилей у тренера, а \texttt{STRING\_AGG} собирает их названия в одну строку через запятую. Используется \texttt{LEFT JOIN}, чтобы в результат попадали и тренеры без специализаций.

\begin{minted}[breaklines]{sql}
SELECT
    t.id,
    t.last_name,
    t.first_name,
    COUNT(tc.car_id) AS car_count,
    STRING_AGG(c.title, ', ') AS cars
FROM trainer t
LEFT JOIN trainer_car tc ON tc.trainer_id = t.id
LEFT JOIN car c ON c.id = tc.car_id
GROUP BY t.id, t.last_name, t.first_name
ORDER BY car_count DESC;
\end{minted}

\begin{table}[H]
\centering
\footnotesize
\setlength{\tabcolsep}{2pt}
\caption{Результат запроса 4}\label{tab:agg4}
\begin{tabular}{|>{\raggedright\arraybackslash}p{8mm}|>{\raggedright\arraybackslash}p{26mm}|>{\raggedright\arraybackslash}p{26mm}|>{\raggedright\arraybackslash}p{22mm}|>{\raggedright\arraybackslash}p{80mm}|}
\hline
\textbf{id} & \textbf{last\_\allowbreak name} & \textbf{first\_\allowbreak name} & \textbf{car\_\allowbreak count} & \textbf{cars} \\ \hline
4 & Иванов & Иван & 2 & ВАЗ 2107, Mitsubishi Lancer Evolution \\ \hline
6 & Иванова & Мария & 2 & Audi RS2, Audi Quattro \\ \hline
2 & Пехота & Наталья & 2 & Audi Quattro, BMW M3 \\ \hline
9 & Петров & Дмитрий & 2 & ВАЗ 2107, BMW M3 \\ \hline
3 & Моралес & Даниэль & 2 & Peugeot 406, ОКА \\ \hline
1 & Торрето & Доминик & 2 & Mazda MX-5 Miata, Toyota Supra \\ \hline
5 & Смирнов & Алексей & 2 & Toyota Supra, Porsche 911 \\ \hline
8 & Соколова & Анна & 2 & Mazda MX-5 Miata, Porsche 911 \\ \hline
\end{tabular}
\end{table}


\textbf{Запрос 5.} Поиск тренеров, которые работают с тремя и более автомобилями. Условие на количество записано в предложении \texttt{HAVING}, так как оно относится к результату агрегатной функции.

\begin{minted}[breaklines]{sql}
SELECT
    t.id,
    t.last_name,
    t.first_name,
    COUNT(tc.car_id) AS car_count
FROM trainer t
JOIN trainer_car tc ON tc.trainer_id = t.id
JOIN car c ON c.id = tc.car_id
WHERE c.is_utilized = FALSE
GROUP BY t.id, t.last_name, t.first_name
HAVING COUNT(tc.car_id) > 2
ORDER BY car_count DESC;
\end{minted}

\begin{table}[H]
\centering
\footnotesize

\caption{Результат запроса 5}\label{tab:agg5}
\begin{tabular}{|l|l|l|l|}
\hline
\textbf{id} & \textbf{last\_\allowbreak name} & \textbf{first\_\allowbreak name} & \textbf{car\_\allowbreak count} \\ \hline
1 & Торрето & Доминик & 3 \\ \hline
\end{tabular}
\end{table}

Условию удовлетворяет один тренер --- Торрето Доминик, который работает с тремя автомобилями.


\subsection{Операторы, встроенные функции и работа с датами}

\textbf{Запрос 6.} Проверка попадания значения в промежуток с помощью оператора \texttt{BETWEEN}: выбираются автомобили мощностью от 400 до 700 л.с., не находящиеся в утилизации.

\begin{minted}[breaklines]{sql}
SELECT *
FROM car
WHERE horsepower BETWEEN 400 AND 700
  AND is_utilized = FALSE;
\end{minted}

\begin{table}[H]
\centering
\scriptsize
\setlength{\tabcolsep}{2pt}
\caption{Результат запроса 6}\label{tab:infun1}
\begin{tabular}{|>{\raggedright\arraybackslash}p{5mm}|>{\raggedright\arraybackslash}p{16mm}|>{\raggedright\arraybackslash}p{21mm}|>{\raggedright\arraybackslash}p{26mm}|>{\raggedright\arraybackslash}p{19mm}|>{\raggedright\arraybackslash}p{19mm}|>{\raggedright\arraybackslash}p{16mm}|>{\raggedright\arraybackslash}p{19mm}|>{\raggedright\arraybackslash}p{14mm}|}
\hline
\textbf{id} & \textbf{title} & \textbf{photo\_\allowbreak url} & \textbf{description} & \textbf{horsepower} & \textbf{production\_\allowbreak year} & \textbf{color} & \textbf{suspension\_\allowbreak name} & \textbf{is\_\allowbreak utilized} \\ \hline
7 & BMW M3 & \url{https://cars.pics/bmw_m3.jpg} & Рядная шестёрка, задний привод. Седан, а едет как спорткар. & 510 & 2021 & Серый & Adaptive M Suspension & false \\ \hline
10 & Porsche 911 & \url{https://cars.pics/porsche_911.jpg} & Оппозитник сзади. Уже 60 лет не меняют форму кузова. Могут, а зачем?. & 450 & 2022 & Чёрный & Sport Suspension & false \\ \hline
\end{tabular}
\end{table}


\textbf{Запрос 7.} Поиск по шаблону с помощью оператора \texttt{LIKE}. Символ \texttt{\%} в шаблоне заменяет любое количество символов, поэтому выбираются все автомобили, название которых начинается с «Audi».

\begin{minted}[breaklines]{sql}
SELECT *
FROM car
WHERE title LIKE 'Audi%';
\end{minted}

\begin{table}[H]
\centering
\scriptsize
\setlength{\tabcolsep}{2pt}
\caption{Результат запроса 7}\label{tab:infun2}
\begin{tabular}{|>{\raggedright\arraybackslash}p{5mm}|>{\raggedright\arraybackslash}p{16mm}|>{\raggedright\arraybackslash}p{21mm}|>{\raggedright\arraybackslash}p{26mm}|>{\raggedright\arraybackslash}p{19mm}|>{\raggedright\arraybackslash}p{19mm}|>{\raggedright\arraybackslash}p{16mm}|>{\raggedright\arraybackslash}p{19mm}|>{\raggedright\arraybackslash}p{14mm}|}
\hline
\textbf{id} & \textbf{title} & \textbf{photo\_\allowbreak url} & \textbf{description} & \textbf{horsepower} & \textbf{production\_\allowbreak year} & \textbf{color} & \textbf{suspension\_\allowbreak name} & \textbf{is\_\allowbreak utilized} \\ \hline
11 & Audi RS2 & \url{https://cars.pics/audi_rs2.jpg} & Универсал, который доводили сами инженеры из Porsche. Первый RS от Audi. & 315 & 1994 & Синий & Sport Suspension & false \\ \hline
8 & Audi Quattro & \url{https://cars.pics/audi_quattro.jpg} & Полный привод и турбонадув. Машина, с которой началась эпоха quattro в ралли. & 200 & 1984 & Оранжевый & MacPherson Strut & false \\ \hline
\end{tabular}
\end{table}


\textbf{Запрос 8.} Формирование полного имени тренера функцией \texttt{CONCAT}. Функция \texttt{COALESCE} возвращает первый аргумент, отличный от \texttt{NULL}, поэтому отсутствующее отчество заменяется пустой строкой. В отличие от оператора \texttt{||}, функция \texttt{CONCAT} воспринимает \texttt{NULL} как пустую строку, и имя не превращается в \texttt{NULL}.

\begin{minted}[breaklines]{sql}
SELECT
    id,
    CONCAT(last_name, ' ', first_name, ' ',
           COALESCE(middle_name, '')) AS full_name
FROM trainer
ORDER BY id;
\end{minted}

\begin{table}[H]
\centering
\footnotesize

\caption{Результат запроса 8}\label{tab:infun3}
\begin{tabular}{|l|l|}
\hline
\textbf{id} & \textbf{full\_\allowbreak name} \\ \hline
1 & Торрето Доминик \\ \hline
2 & Пехота Наталья Морская \\ \hline
3 & Моралес Даниэль \\ \hline
4 & Иванов Иван \\ \hline
5 & Смирнов Алексей \\ \hline
6 & Иванова Мария \\ \hline
8 & Соколова Анна \\ \hline
9 & Петров Дмитрий \\ \hline
\end{tabular}
\end{table}


\textbf{Запрос 9.} Конструкция \texttt{CASE} используется для классификации автомобилей по мощности на три категории.

\begin{minted}[breaklines]{sql}
SELECT
    title,
    horsepower,
    CASE
        WHEN horsepower < 300 THEN 'Низкая мощность'
        WHEN horsepower BETWEEN 300 AND 600 THEN 'Средняя мощность'
        ELSE 'Высокая мощность'
    END AS power_category
FROM car
WHERE is_utilized = FALSE
ORDER BY horsepower DESC;
\end{minted}

\begin{table}[H]
\centering
\footnotesize

\caption{Результат запроса 9}\label{tab:infun4}
\begin{tabular}{|l|l|l|}
\hline
\textbf{title} & \textbf{horsepower} & \textbf{power\_\allowbreak category} \\ \hline
BMW M3 & 510 & Средняя мощность \\ \hline
Porsche 911 & 450 & Средняя мощность \\ \hline
Toyota Supra & 330 & Средняя мощность \\ \hline
Audi RS2 & 315 & Средняя мощность \\ \hline
Nissan Skyline GT-R & 280 & Низкая мощность \\ \hline
Mazda MX-5 Miata & 201 & Низкая мощность \\ \hline
Audi Quattro & 200 & Низкая мощность \\ \hline
Peugeot 406 & 194 & Низкая мощность \\ \hline
Mitsubishi Lancer Evolution & 180 & Низкая мощность \\ \hline
Volkswagen Beetle & 111 & Низкая мощность \\ \hline
ВАЗ 2107 & 76 & Низкая мощность \\ \hline
\end{tabular}
\end{table}


\textbf{Запрос 10.} Вычисление примерного возраста тренеров. Функция \texttt{EXTRACT} получает год из текущей даты, результат приводится к целому типу оператором \texttt{::INT}, после чего из него вычитается год рождения.

\begin{minted}[breaklines]{sql}
SELECT
    id,
    last_name,
    first_name,
    birth_year,
    EXTRACT(YEAR FROM CURRENT_DATE)::INT - birth_year AS approximate_age
FROM trainer
ORDER BY birth_year;
\end{minted}

\begin{table}[H]
\centering
\footnotesize

\caption{Результат запроса 10}\label{tab:infun5}
\begin{tabular}{|l|l|l|l|l|}
\hline
\textbf{id} & \textbf{last\_\allowbreak name} & \textbf{first\_\allowbreak name} & \textbf{birth\_\allowbreak year} & \textbf{approximate\_\allowbreak age} \\ \hline
1 & Торрето & Доминик & 1967 & 59 \\ \hline
3 & Моралес & Даниэль & 1973 & 53 \\ \hline
9 & Петров & Дмитрий & 1979 & 47 \\ \hline
4 & Иванов & Иван & 1985 & 41 \\ \hline
2 & Пехота & Наталья & 1988 & 38 \\ \hline
6 & Иванова & Мария & 1990 & 36 \\ \hline
5 & Смирнов & Алексей & 1992 & 34 \\ \hline
8 & Соколова & Анна & 1995 & 31 \\ \hline
\end{tabular}
\end{table}


\textbf{Запрос 11.} Работа с временными интервалами. Из текущей даты вычитается интервал в пять лет, и выбираются соревнования, прошедшие в этот промежуток.

\begin{minted}[breaklines]{sql}
SELECT
    id,
    name,
    event_date
FROM competition
WHERE event_date BETWEEN CURRENT_DATE - INTERVAL '5 years' AND CURRENT_DATE
ORDER BY event_date DESC;
\end{minted}

\begin{table}[H]
\centering
\footnotesize

\caption{Результат запроса 11}\label{tab:infun6}
\begin{tabular}{|l|l|l|}
\hline
\textbf{id} & \textbf{name} & \textbf{event\_\allowbreak date} \\ \hline
11 & Кубок Поршня & 2023-06-15 \\ \hline
12 & Гонка Чемпионов & 2022-09-03 \\ \hline
\end{tabular}
\end{table}


\textbf{Запрос 12.} Вычитание дат. Разность двух значений типа \texttt{date} равна количеству дней между ними.

\begin{minted}[breaklines]{sql}
SELECT
    name,
    event_date,
    CURRENT_DATE - event_date AS days_since_competition
FROM competition
WHERE event_date <= CURRENT_DATE
ORDER BY event_date DESC;
\end{minted}

\begin{table}[H]
\centering
\footnotesize

\caption{Результат запроса 12}\label{tab:infun7}
\begin{tabular}{|l|l|l|}
\hline
\textbf{name} & \textbf{event\_\allowbreak date} & \textbf{days\_\allowbreak since\_\allowbreak competition} \\ \hline
Кубок Поршня & 2023-06-15 & 1230 \\ \hline
Гонка Чемпионов & 2022-09-03 & 1515 \\ \hline
Ралли Монте-Карло & 2021-01-24 & 2102 \\ \hline
Ралли Финляндия & 2020-08-07 & 2272 \\ \hline
Ралли Греция & 2019-06-02 & 2704 \\ \hline
Ралли Уэльс & 2018-10-05 & 2944 \\ \hline
Ралли Мичиган & 2017-07-14 & 3392 \\ \hline
Ралли Испания & 2016-10-16 & 3663 \\ \hline
Ралли Марокко & 2015-04-26 & 4202 \\ \hline
Ралли Австралия & 2014-09-12 & 4428 \\ \hline
\end{tabular}
\end{table}


\subsection{Подзапросы и предложение WITH}

Результаты одного запроса можно использовать в другом. Подзапрос может возвращать одно значение (скалярный подзапрос), набор значений или использоваться для проверки существования строк. Общие табличные выражения (CTE), задаваемые в предложении \texttt{WITH}, позволяют дать подзапросу имя, вычислить его один раз и использовать результат в основном запросе.

\textbf{Запрос 13.} Скалярный подзапрос в условии \texttt{WHERE}: сначала вычисляется средняя мощность автомобилей автопарка, затем выбираются автомобили с большей мощностью.

\begin{minted}[breaklines]{sql}
SELECT *
FROM car
WHERE horsepower > (SELECT AVG(horsepower) FROM car);
\end{minted}

\begin{table}[H]
\centering
\scriptsize
\setlength{\tabcolsep}{2pt}
\caption{Результат запроса 13}\label{tab:cte1}
\begin{tabular}{|>{\raggedright\arraybackslash}p{5mm}|>{\raggedright\arraybackslash}p{16mm}|>{\raggedright\arraybackslash}p{21mm}|>{\raggedright\arraybackslash}p{26mm}|>{\raggedright\arraybackslash}p{19mm}|>{\raggedright\arraybackslash}p{19mm}|>{\raggedright\arraybackslash}p{16mm}|>{\raggedright\arraybackslash}p{19mm}|>{\raggedright\arraybackslash}p{14mm}|}
\hline
\textbf{id} & \textbf{title} & \textbf{photo\_\allowbreak url} & \textbf{description} & \textbf{horsepower} & \textbf{production\_\allowbreak year} & \textbf{color} & \textbf{suspension\_\allowbreak name} & \textbf{is\_\allowbreak utilized} \\ \hline
4 & Toyota Supra & \url{https://cars.pics/toyota_supra.jpg} & Мотор 2JZ, тюнинг без границ. Такую вы видели только на картинках! & 330 & 1997 & Оранжевый & Double Wishbone & false \\ \hline
5 & Nissan Skyline GT-R & \url{https://cars.pics/nissan_skyline_gtr.jpg} & Его знает каждый. Столько игр и фильмов ему посвящено! & 280 & 1999 & Синий & Multi-Link & false \\ \hline
7 & BMW M3 & \url{https://cars.pics/bmw_m3.jpg} & Рядная шестёрка, задний привод. Седан, а едет как спорткар. & 510 & 2021 & Серый & Adaptive M Suspension & false \\ \hline
10 & Porsche 911 & \url{https://cars.pics/porsche_911.jpg} & Оппозитник сзади. Уже 60 лет не меняют форму кузова. Могут, а зачем?. & 450 & 2022 & Чёрный & Sport Suspension & false \\ \hline
11 & Audi RS2 & \url{https://cars.pics/audi_rs2.jpg} & Универсал, который доводили сами инженеры из Porsche. Первый RS от Audi. & 315 & 1994 & Синий & Sport Suspension & false \\ \hline
\end{tabular}
\end{table}


\textbf{Запрос 14.} Коррелированный подзапрос с предикатом \texttt{NOT EXISTS}: выбираются тренеры, для которых в таблице \texttt{achievement} нет ни одной записи.

\begin{minted}[breaklines]{sql}
SELECT
    t.id,
    t.last_name,
    t.first_name
FROM trainer t
WHERE NOT EXISTS (
    SELECT 1
    FROM achievement a
    WHERE a.trainer_id = t.id
);
\end{minted}

\begin{table}[H]
\centering
\footnotesize

\caption{Результат запроса 14}\label{tab:cte2}
\begin{tabular}{|l|l|l|}
\hline
\textbf{id} & \textbf{last\_\allowbreak name} & \textbf{first\_\allowbreak name} \\ \hline
8 & Соколова & Анна \\ \hline
\end{tabular}
\end{table}

Не имеет достижений в соревнованиях один тренер --- Соколова Анна.


\textbf{Запрос 15.} Общее табличное выражение \texttt{trainer\_results} подсчитывает количество достижений каждого тренера. Далее результат соединяется с таблицей \texttt{trainer} для получения фамилий и имён.

\begin{minted}[breaklines]{sql}
WITH trainer_results AS (
    SELECT
        trainer_id,
        COUNT(*) AS achievement_count
    FROM achievement
    GROUP BY trainer_id
)
SELECT
    t.last_name,
    t.first_name,
    tr.achievement_count
FROM trainer_results tr
INNER JOIN trainer t ON t.id = tr.trainer_id
ORDER BY tr.achievement_count DESC;
\end{minted}

\begin{table}[H]
\centering
\footnotesize

\caption{Результат запроса 15}\label{tab:cte3}
\begin{tabular}{|l|l|l|}
\hline
\textbf{last\_\allowbreak name} & \textbf{first\_\allowbreak name} & \textbf{achievement\_\allowbreak count} \\ \hline
Моралес & Даниэль & 2 \\ \hline
Пехота & Наталья & 2 \\ \hline
Торрето & Доминик & 2 \\ \hline
Иванова & Мария & 1 \\ \hline
Иванов & Иван & 1 \\ \hline
Смирнов & Алексей & 1 \\ \hline
Петров & Дмитрий & 1 \\ \hline
\end{tabular}
\end{table}


\subsection{Предложение RETURNING}

Команды \texttt{INSERT}, \texttt{UPDATE} и \texttt{DELETE} могут возвращать затронутые строки с помощью предложения \texttt{RETURNING}. Для \texttt{UPDATE} возвращаются уже изменённые значения.

\textbf{Запрос 16.} Добавление нового автомобиля. Предложение \texttt{RETURNING *} возвращает созданную запись целиком, в том числе автоматически сгенерированный идентификатор и значения по умолчанию.

\begin{minted}[breaklines]{sql}
INSERT INTO car (
    title,
    description,
    horsepower,
    production_year,
    color,
    suspension_name
)
VALUES (
    'Volkswagen Beetle',
    'Самый стильный автомобиль в клубе',
    111,
    1975,
    'Красный',
    'MacPherson'
)
RETURNING *;
\end{minted}

\begin{table}[H]
\centering
\scriptsize
\setlength{\tabcolsep}{2pt}
\caption{Результат запроса 16}\label{tab:ret1}
\begin{tabular}{|>{\raggedright\arraybackslash}p{5mm}|>{\raggedright\arraybackslash}p{16mm}|>{\raggedright\arraybackslash}p{21mm}|>{\raggedright\arraybackslash}p{26mm}|>{\raggedright\arraybackslash}p{19mm}|>{\raggedright\arraybackslash}p{19mm}|>{\raggedright\arraybackslash}p{16mm}|>{\raggedright\arraybackslash}p{19mm}|>{\raggedright\arraybackslash}p{14mm}|}
\hline
\textbf{id} & \textbf{title} & \textbf{photo\_\allowbreak url} & \textbf{description} & \textbf{horsepower} & \textbf{production\_\allowbreak year} & \textbf{color} & \textbf{suspension\_\allowbreak name} & \textbf{is\_\allowbreak utilized} \\ \hline
15 & Volkswagen Beetle & null & Самый стильный автомобиль в клубе & 111 & 1975 & Красный & MacPherson & false \\ \hline
\end{tabular}
\end{table}


\textbf{Запрос 17.} Увеличение мощности автомобиля на 20 л.с. Возвращаются только изменённые значения выбранных столбцов.

\begin{minted}[breaklines]{sql}
UPDATE car
SET horsepower = horsepower + 20
WHERE id = 3
RETURNING id, title, horsepower;
\end{minted}

\begin{table}[H]
\centering
\footnotesize

\caption{Результат запроса 17}\label{tab:ret2}
\begin{tabular}{|l|l|l|}
\hline
\textbf{id} & \textbf{title} & \textbf{horsepower} \\ \hline
3 & Mazda MX-5 Miata & 221 \\ \hline
\end{tabular}
\end{table}


\textbf{Запрос 18.} Удаление специализации тренера Торрето на автомобиле Nissan Skyline GT-R. Идентификаторы тренера и автомобиля определяются подзапросами, удалённая строка возвращается в результате.

\begin{minted}[breaklines]{sql}
DELETE FROM trainer_car
WHERE trainer_id IN (SELECT id FROM trainer WHERE last_name = 'Торрето')
    AND car_id IN (SELECT id FROM car WHERE title = 'Nissan Skyline GT-R')
RETURNING *;
\end{minted}

\begin{table}[H]
\centering
\footnotesize

\caption{Результат запроса 18}\label{tab:ret3}
\begin{tabular}{|l|l|}
\hline
\textbf{trainer\_\allowbreak id} & \textbf{car\_\allowbreak id} \\ \hline
1 & 5 \\ \hline
\end{tabular}
\end{table}


\subsection{Планы выполнения запросов (EXPLAIN)}

План выполнения запроса представляет собой дерево узлов, для каждого из которых выводятся оценка стоимости запуска и общая стоимость (\texttt{cost}), ожидаемое число строк (\texttt{rows}) и ширина строки (\texttt{width}). При использовании параметра \texttt{ANALYZE} запрос реально выполняется, и к оценкам добавляются фактическое время и число строк.

\textbf{План 1.} Обычный \texttt{EXPLAIN} для поиска автомобиля по названию.

\begin{minted}[breaklines]{sql}
EXPLAIN
SELECT *
FROM car
WHERE title = 'Mazda MX-5 Miata';
\end{minted}

\begin{Verbatim}[fontsize=\scriptsize,breaklines=true,breakanywhere=true]
Index Scan using car_title_idx on car  (cost=0.14..8.16 rows=1 width=1059)
Index Cond: ((title)::text = 'Mazda MX-5 Miata'::text)
\end{Verbatim}


По столбцу \texttt{title} создан индекс, поэтому планировщик выбрал метод доступа \textit{Index Scan} с условием поиска по индексу \texttt{car\_title\_idx}. Общая оценка стоимости составляет 8.16.


\textbf{План 2.} \texttt{EXPLAIN (ANALYZE)} для запроса с \texttt{NOT EXISTS}. Помимо оценок выводится фактическое время выполнения.

\begin{minted}[breaklines]{sql}
EXPLAIN (ANALYSE)
SELECT
    t.id,
    t.last_name,
    t.first_name
FROM trainer t
WHERE NOT EXISTS (
    SELECT 1
    FROM achievement a
    WHERE a.trainer_id = t.id
);
\end{minted}

\begin{Verbatim}[fontsize=\scriptsize,breaklines=true,breakanywhere=true]
Nested Loop Anti Join  (cost=0.15..32.62 rows=30 width=444) (actual time=0.032..0.034 rows=1.00 loops=1)
Buffers: shared hit=16
->  Seq Scan on trainer t  (cost=0.00..10.60 rows=60 width=444) (actual time=0.014..0.015 rows=8.00 loops=1)
Buffers: shared hit=1
->  Index Only Scan using achievement_trainer_id_idx on achievement a  (cost=0.15..1.02 rows=4 width=8) (actual time=0.002..0.002 rows=0.88 loops=8)
Index Cond: (trainer_id = t.id)
Heap Fetches: 7
Index Searches: 8
Buffers: shared hit=15
Planning Time: 0.111 ms
Execution Time: 0.048 ms
\end{Verbatim}


Для соединения использован узел \textit{Nested Loop Anti Join}: таблица \texttt{trainer} читается последовательно (\textit{Seq Scan}), а для каждого тренера выполняется поиск в индексе \texttt{achievement\_trainer\_id\_idx} (\textit{Index Only Scan}). Время планирования составило 0.111 мс, время выполнения --- 0.048 мс. Оценка числа строк таблицы \texttt{trainer} (60) заметно отличается от фактической (8), что объясняется малым объёмом данных и отсутствием собранной статистики.


\textbf{План 3.} \texttt{EXPLAIN (ANALYZE, BUFFERS)} с получением только столбца \texttt{title}. Поскольку этот столбец входит в индекс, применяется метод доступа \textit{Index Only Scan}.

\begin{minted}[breaklines]{sql}
EXPLAIN (ANALYZE, BUFFERS)
SELECT title
FROM car
WHERE title = 'Toyota Supra';
\end{minted}

\begin{Verbatim}[fontsize=\scriptsize,breaklines=true,breakanywhere=true]
Index Only Scan using car_title_idx on car  (cost=0.14..8.16 rows=1 width=218) (actual time=0.022..0.022 rows=1.00 loops=1)
Index Cond: (title = 'Toyota Supra'::text)
Heap Fetches: 1
Index Searches: 1
Buffers: shared hit=2
Planning Time: 0.054 ms
Execution Time: 0.032 ms
\end{Verbatim}


Параметр \texttt{BUFFERS} добавил информацию об использовании буферов: \texttt{shared hit=2} означает, что оба нужных блока найдены в кеше, и чтения с диска не потребовалось. Значение \texttt{Heap Fetches: 1} показывает, что для найденной строки пришлось обратиться к таблице, так как страница ещё не отмечена в карте видимости (таблица недавно изменялась).


\subsection{План запроса в формате JSON для нескольких связанных таблиц}

\textbf{План 4.} Формат \texttt{JSON} содержит ту же информацию, что и текстовый, но удобнее для программной обработки. Выборка объединяет три связанные таблицы: \texttt{trainer}, \texttt{trainer\_car} и \texttt{car}.

\begin{minted}[breaklines]{sql}
EXPLAIN (ANALYSE, FORMAT JSON)
SELECT
    t.last_name,
    t.first_name,
    c.title AS car_title
FROM trainer t
JOIN trainer_car tc ON tc.trainer_id = t.id
JOIN car c ON c.id = tc.car_id;
\end{minted}

Результат выполнения:

\begin{Verbatim}[fontsize=\tiny,breaklines=true,breakanywhere=true]
[
  {
    "Plan": {
      "Node Type": "Hash Join",
      "Parallel Aware": false,
      "Async Capable": false,
      "Join Type": "Inner",
      "Startup Cost": 22.92,
      "Total Cost": 61.34,
      "Plan Rows": 1850,
      "Plan Width": 654,
      "Actual Startup Time": 0.045,
      "Actual Total Time": 0.051,
      "Actual Rows": 17.00,
      "Actual Loops": 1,
      "Disabled": false,
      "Inner Unique": true,
      "Hash Cond": "(tc.car_id = c.id)",
      "Shared Hit Blocks": 3,
      "Shared Read Blocks": 0,
      "Shared Dirtied Blocks": 0,
      "Shared Written Blocks": 0,
      "Local Hit Blocks": 0,
      "Local Read Blocks": 0,
      "Local Dirtied Blocks": 0,
      "Local Written Blocks": 0,
      "Temp Read Blocks": 0,
      "Temp Written Blocks": 0,
      "Plans": [
        {
          "Node Type": "Hash Join",
          "Parent Relationship": "Outer",
          "Parallel Aware": false,
          "Async Capable": false,
          "Join Type": "Inner",
          "Startup Cost": 11.35,
          "Total Cost": 44.81,
          "Plan Rows": 1850,
          "Plan Width": 444,
          "Actual Startup Time": 0.019,
          "Actual Total Time": 0.022,
          "Actual Rows": 17.00,
          "Actual Loops": 1,
          "Disabled": false,
          "Inner Unique": true,
          "Hash Cond": "(tc.trainer_id = t.id)",
          "Shared Hit Blocks": 2,
          "Shared Read Blocks": 0,
          "Shared Dirtied Blocks": 0,
          "Shared Written Blocks": 0,
          "Local Hit Blocks": 0,
          "Local Read Blocks": 0,
          "Local Dirtied Blocks": 0,
          "Local Written Blocks": 0,
          "Temp Read Blocks": 0,
          "Temp Written Blocks": 0,
          "Plans": [
            {
              "Node Type": "Seq Scan",
              "Parent Relationship": "Outer",
              "Parallel Aware": false,
              "Async Capable": false,
              "Relation Name": "trainer_car",
              "Alias": "tc",
              "Startup Cost": 0.00,
              "Total Cost": 28.50,
              "Plan Rows": 1850,
              "Plan Width": 16,
              "Actual Startup Time": 0.005,
              "Actual Total Time": 0.006,
              "Actual Rows": 17.00,
              "Actual Loops": 1,
              "Disabled": false,
              "Shared Hit Blocks": 1,
              "Shared Read Blocks": 0,
              "Shared Dirtied Blocks": 0,
              "Shared Written Blocks": 0,
              "Local Hit Blocks": 0,
              "Local Read Blocks": 0,
              "Local Dirtied Blocks": 0,
              "Local Written Blocks": 0,
              "Temp Read Blocks": 0,
              "Temp Written Blocks": 0
            },
            {
              "Node Type": "Hash",
              "Parent Relationship": "Inner",
              "Parallel Aware": false,
              "Async Capable": false,
              "Startup Cost": 10.60,
              "Total Cost": 10.60,
              "Plan Rows": 60,
              "Plan Width": 444,
              "Actual Startup Time": 0.008,
              "Actual Total Time": 0.008,
              "Actual Rows": 8.00,
              "Actual Loops": 1,
              "Disabled": false,
              "Hash Buckets": 1024,
              "Original Hash Buckets": 1024,
              "Hash Batches": 1,
              "Original Hash Batches": 1,
              "Peak Memory Usage": 9,
              "Shared Hit Blocks": 1,
              "Shared Read Blocks": 0,
              "Shared Dirtied Blocks": 0,
              "Shared Written Blocks": 0,
              "Local Hit Blocks": 0,
              "Local Read Blocks": 0,
              "Local Dirtied Blocks": 0,
              "Local Written Blocks": 0,
              "Temp Read Blocks": 0,
              "Temp Written Blocks": 0,
              "Plans": [
                {
                  "Node Type": "Seq Scan",
                  "Parent Relationship": "Outer",
                  "Parallel Aware": false,
                  "Async Capable": false,
                  "Relation Name": "trainer",
                  "Alias": "t",
                  "Startup Cost": 0.00,
                  "Total Cost": 10.60,
                  "Plan Rows": 60,
                  "Plan Width": 444,
                  "Actual Startup Time": 0.005,
                  "Actual Total Time": 0.006,
                  "Actual Rows": 8.00,
                  "Actual Loops": 1,
                  "Disabled": false,
                  "Shared Hit Blocks": 1,
                  "Shared Read Blocks": 0,
                  "Shared Dirtied Blocks": 0,
                  "Shared Written Blocks": 0,
                  "Local Hit Blocks": 0,
                  "Local Read Blocks": 0,
                  "Local Dirtied Blocks": 0,
                  "Local Written Blocks": 0,
                  "Temp Read Blocks": 0,
                  "Temp Written Blocks": 0
                }
              ]
            }
          ]
        },
        {
          "Node Type": "Hash",
          "Parent Relationship": "Inner",
          "Parallel Aware": false,
          "Async Capable": false,
          "Startup Cost": 10.70,
          "Total Cost": 10.70,
          "Plan Rows": 70,
          "Plan Width": 226,
          "Actual Startup Time": 0.019,
          "Actual Total Time": 0.020,
          "Actual Rows": 12.00,
          "Actual Loops": 1,
          "Disabled": false,
          "Hash Buckets": 1024,
          "Original Hash Buckets": 1024,
          "Hash Batches": 1,
          "Original Hash Batches": 1,
          "Peak Memory Usage": 9,
          "Shared Hit Blocks": 1,
          "Shared Read Blocks": 0,
          "Shared Dirtied Blocks": 0,
          "Shared Written Blocks": 0,
          "Local Hit Blocks": 0,
          "Local Read Blocks": 0,
          "Local Dirtied Blocks": 0,
          "Local Written Blocks": 0,
          "Temp Read Blocks": 0,
          "Temp Written Blocks": 0,
          "Plans": [
            {
              "Node Type": "Seq Scan",
              "Parent Relationship": "Outer",
              "Parallel Aware": false,
              "Async Capable": false,
              "Relation Name": "car",
              "Alias": "c",
              "Startup Cost": 0.00,
              "Total Cost": 10.70,
              "Plan Rows": 70,
              "Plan Width": 226,
              "Actual Startup Time": 0.013,
              "Actual Total Time": 0.015,
              "Actual Rows": 12.00,
              "Actual Loops": 1,
              "Disabled": false,
              "Shared Hit Blocks": 1,
              "Shared Read Blocks": 0,
              "Shared Dirtied Blocks": 0,
              "Shared Written Blocks": 0,
              "Local Hit Blocks": 0,
              "Local Read Blocks": 0,
              "Local Dirtied Blocks": 0,
              "Local Written Blocks": 0,
              "Temp Read Blocks": 0,
              "Temp Written Blocks": 0
            }
          ]
        }
      ]
    },
    "Planning": {
      "Shared Hit Blocks": 0,
      "Shared Read Blocks": 0,
      "Shared Dirtied Blocks": 0,
      "Shared Written Blocks": 0,
      "Local Hit Blocks": 0,
      "Local Read Blocks": 0,
      "Local Dirtied Blocks": 0,
      "Local Written Blocks": 0,
      "Temp Read Blocks": 0,
      "Temp Written Blocks": 0
    },
    "Planning Time": 0.162,
    "Triggers": [
    ],
    "Execution Time": 0.069
  }
]
\end{Verbatim}


Корневой узел плана --- \textit{Hash Join}, общая стоимость 61.34, фактическое время 0.051 мс. Внутри него расположен второй \textit{Hash Join}, соединяющий \texttt{trainer\_car} и \texttt{trainer}; все три таблицы читаются последовательно (\textit{Seq Scan}). В результате получено 17 строк, время планирования составило 0.162 мс, время выполнения --- 0.069 мс.


\subsection{Способы соединения наборов строк}

Для одной и той же выборки (тренеры и автомобили, на которых они специализируются) выполнено соединение тремя способами: вложенным циклом (\textit{Nested Loop}), хешированием (\textit{Hash Join}) и слиянием (\textit{Merge Join}). Нужный способ включается принудительно: параметрами \texttt{enable\_\allowbreak nestloop}, \texttt{enable\_\allowbreak hashjoin} и \texttt{enable\_\allowbreak mergejoin} отключаются остальные способы. Для каждого варианта получен план в формате \texttt{JSON} с параметром \texttt{ANALYZE}. Ниже приведено сокращённое представление плана (тип узла, стоимость, фактическое время и число строк) и основные показатели.

\textbf{Вариант 1. Вложенный цикл.}

\begin{minted}[breaklines]{sql}
SET enable_nestloop = ON;
SET enable_hashjoin = OFF;
SET enable_mergejoin = OFF;

EXPLAIN (ANALYZE, FORMAT JSON)
SELECT
    t.id,
    t.last_name,
    c.title
FROM trainer t
INNER JOIN trainer_car tc ON tc.trainer_id = t.id
INNER JOIN car c ON c.id = tc.car_id;
\end{minted}

\begin{Verbatim}[fontsize=\footnotesize,breaklines=true,breakanywhere=true]
Nested Loop  (cost=0.29..425.70, actual time=0.037..0.077, rows=17)
   -> Nested Loop  (cost=0.15..86.70, actual time=0.030..0.057, rows=17)
      -> Seq Scan on car  (cost=0.00..10.70, actual time=0.014..0.015, rows=12)
      -> Index Scan on trainer_car using trainer_car_car_id_idx  (cost=0.15..1.00, actual time=0.002..0.003, rows=1)
   -> Index Scan on trainer using trainer_pkey  (cost=0.14..0.18, actual time=0.001..0.001, rows=1)
\end{Verbatim}


Из таблицы \texttt{car} (12 строк) выполняется последовательное чтение, а для каждой строки поиск соответствующих записей в \texttt{trainer\_car} и \texttt{trainer} выполняется по индексам \texttt{trainer\_car\_car\_id\_idx} и \texttt{trainer\_pkey}. Большое число обращений к индексам даёт самую высокую стоимость плана.


\textbf{Вариант 2. Хеширование.}

\begin{minted}[breaklines]{sql}
SET enable_nestloop = OFF;
SET enable_hashjoin = ON;
SET enable_mergejoin = OFF;

EXPLAIN (ANALYZE, FORMAT JSON)
SELECT
    t.id,
    t.last_name,
    c.title
FROM trainer t
JOIN trainer_car tc ON tc.trainer_id = t.id
JOIN car c ON c.id = tc.car_id;
\end{minted}

\begin{Verbatim}[fontsize=\footnotesize,breaklines=true,breakanywhere=true]
Hash Join  (cost=22.92..61.34, actual time=0.053..0.061, rows=17)
   -> Hash Join  (cost=11.35..44.81, actual time=0.022..0.026, rows=17)
      -> Seq Scan on trainer_car  (cost=0.00..28.50, actual time=0.006..0.007, rows=17)
      -> Hash  (cost=10.60..10.60, actual time=0.010..0.010, rows=8)
         -> Seq Scan on trainer  (cost=0.00..10.60, actual time=0.006..0.007, rows=8)
   -> Hash  (cost=10.70..10.70, actual time=0.023..0.024, rows=12)
      -> Seq Scan on car  (cost=0.00..10.70, actual time=0.017..0.018, rows=12)
\end{Verbatim}


Меньшие наборы строк помещаются в хеш-таблицы, после чего большая таблица \texttt{trainer\_car} читается один раз последовательно. Подготовка хеш-таблиц дешёвая, так как таблицы небольшие.


\textbf{Вариант 3. Слияние.}

\begin{minted}[breaklines]{sql}
SET enable_nestloop = OFF;
SET enable_hashjoin = OFF;
SET enable_mergejoin = ON;

EXPLAIN (ANALYZE, FORMAT JSON)
SELECT
    t.id,
    t.last_name,
    c.title
FROM trainer t
JOIN trainer_car tc ON tc.trainer_id = t.id
JOIN car c ON c.id = tc.car_id;
\end{minted}

\begin{Verbatim}[fontsize=\footnotesize,breaklines=true,breakanywhere=true]
Merge Join  (cost=211.99..228.07, actual time=0.051..0.055, rows=17)
   -> Sort  (cost=199.14..203.77, actual time=0.041..0.042, rows=17)
      -> Merge Join  (cost=12.52..98.75, actual time=0.030..0.037, rows=17)
         -> Index Only Scan on trainer_car using trainer_car_pkey  (cost=0.15..75.90, actual time=0.011..0.015, rows=17)
         -> Sort  (cost=12.37..12.52, actual time=0.016..0.016, rows=8)
            -> Seq Scan on trainer  (cost=0.00..10.60, actual time=0.010..0.011, rows=8)
   -> Sort  (cost=12.85..13.02, actual time=0.009..0.009, rows=11)
      -> Seq Scan on car  (cost=0.00..10.70, actual time=0.005..0.005, rows=12)
\end{Verbatim}


Для слияния оба набора должны быть отсортированы по ключу соединения, поэтому в плане появились узлы \textit{Sort}, а общая стоимость выросла за счёт сортировки. Для таблицы \texttt{trainer\_car} использован \textit{Index Only Scan} по первичному ключу.


Показатели трёх вариантов сведены в таблицу~\ref{tab:joins}: \textit{Actual Time} --- фактическое время корневого узла, \textit{Planning} и \textit{Execution} --- время планирования и выполнения запроса, \textit{Shared Hit} --- число блоков, найденных в разделяемом буфере.

\begin{table}[H]\centering\footnotesize\caption{Сравнение способов соединения}\label{tab:joins}
\begin{tabular}{|l|r|r|r|r|r|}\hline
\textbf{Способ} & \textbf{Total Cost} & \textbf{Actual Time, мс} & \textbf{Planning, мс} & \textbf{Execution, мс} & \textbf{Shared Hit} \\ \hline
Nested Loop & 425.70 & 0.077 & 0.181 & 0.101 & 59 \\ \hline
Hash Join & 61.34 & 0.061 & 0.181 & 0.080 & 3 \\ \hline
Merge Join & 228.07 & 0.055 & 0.159 & 0.079 & 4 \\ \hline
\end{tabular}\end{table}


Результаты показывают, что наименьшую оценку стоимости (61.34) имеет соединение хешированием, самую высокую (425.70) --- вложенный цикл. Время выполнения всех трёх вариантов составляет порядка 0.08--0.10 мс; на выборке из 17 строк разница между ними сопоставима с погрешностью измерения, при этом вложенный цикл оказался самым медленным и обратился к наибольшему числу буферов (59 против 3 у хеширования и 4 у слияния). Планировщик при включённых по умолчанию способах соединения сам выбрал \textit{Hash Join} (см. план 4). Таким образом, для данной выборки наиболее оптимальным способом является \textbf{соединение хешированием (Hash Join)}: оно имеет наименьшую стоимость и наименьшее число обращений к буферам, а время выполнения практически не отличается от лучшего (0.080 мс у хеширования и 0.079 мс у слияния). При росте объёмов данных это преимущество будет увеличиваться, так как хеширование читает каждую таблицу один раз, а вложенный цикл выполняет обращение к индексу для каждой строки внешнего набора.

\section*{Вывод}

В ходе выполнения лабораторной работы были освоены агрегатные функции с группировкой строк и фильтрацией групп (\texttt{GROUP BY}, \texttt{HAVING}), операторы и встроенные функции PostgreSQL, работа с датами и временными интервалами, подзапросы, общие табличные выражения и предложение \texttt{RETURNING}. Все примеры из методических указаний выполнены для базы данных гоночного автоклуба.

Получены и проанализированы планы выполнения запросов с параметрами \texttt{ANALYZE} и \texttt{BUFFERS}, а также в формате \texttt{JSON} для выборки из нескольких связанных таблиц. Рассмотрены методы доступа \textit{Index Scan} и \textit{Index Only Scan}. Три способа соединения (\textit{Nested Loop}, \textit{Hash Join}, \textit{Merge Join}) сравнены на одной выборке; наиболее оптимальным признано соединение хешированием.

\end{document}
